\chapter{MÉTHODES ET TECHNOLOGIES PRÉVUES}
\thispagestyle{empty}

\section{Méthodes}

\textbf{Construction du jeu de données -- Custom-WikiCS :}
Le jeu de données utilisé dans le projet a été reconstruit en se basant sur la topologie du jeu de données de référence Wiki-CS \cite{mernyei2020wikics}. Le processus se compose des étapes suivantes :
\begin{enumerate}
    \item La topologie de graphe du jeu de données Wiki-CS (11 701 n\oe uds, 291 039 arêtes orientées) a été extraite.
    \item Les titres des n\oe uds ont été traités et mis en correspondance avec les titres réels des articles sur Wikipédia.
    \item Les textes des articles correspondants ont été extraits à l'aide du jeu de données ``Wikipedia STEM Corpus'' disponible sur Kaggle.
    \item Les textes des articles absents du corpus ont été complétés via l'API Wikipédia.
    \item Les textes obtenus ont été vectorisés avec deux modèles de plongements de phrases différents pour créer les attributs de n\oe uds.
\end{enumerate}

Cette approche remplace la simple représentation textuelle par moyenne de GloVe du Wiki-CS original par des vecteurs de plongements sémantiques modernes, tout en préservant la fiabilité de la structure de graphe.

\textbf{Modèles de plongements de phrases :}
Deux modèles pré-entraînés différents sont évalués de manière comparative pour créer la représentation vectorielle des textes des articles :
\begin{itemize}
    \item \textbf{BAAI/bge-m3 :} un modèle multilingue et multitâche produisant des vecteurs de plongement de 1024 dimensions.
    \item \textbf{Alibaba-NLP/gte-modernbert-base :} un modèle basé sur l'architecture ModernBERT produisant des vecteurs de plongement de 768 dimensions.
\end{itemize}
Des expériences sont en cours pour déterminer quel modèle fournit les meilleures performances de recommandation lorsqu'il est combiné avec le GNN.

\textbf{Réseau de neurones sur graphes -- GAT :}
Les Graph Attention Networks (GAT) \cite{velickovic2018graph} seront utilisés pour apprendre les relations structurelles entre les articles. Le GAT est capable d'apprendre quel lien est le plus important en attribuant des poids d'attention différents à chaque n\oe ud voisin. Le GCN \cite{kipf2017semi} sera utilisé comme modèle de référence (baseline) pour la comparaison.

\textbf{Grand modèle de langage -- Présentation hiérarchique :}
L'utilisation d'un grand modèle de langage open-source exécutable localement (par ex. DeepSeek-R1 ou un modèle similaire) est prévue pour présenter les résultats de recommandation basés sur le GNN à l'utilisateur sous forme de parcours d'apprentissage hiérarchique. La compatibilité avec une limite de 8 Go de VRAM est prise en compte dans le choix du modèle.

\section{Logiciels}

\begin{itemize}
    \item \textbf{Python :} langage de programmation principal.
    \item \textbf{PyTorch \& PyTorch Geometric (PyG) :} pour la construction et l'entraînement des modèles GNN (GAT, GCN).
    \item \textbf{HuggingFace Transformers / Sentence-Transformers :} pour l'exécution des modèles de plongements de phrases.
    \item \textbf{NetworkX / iGraph :} pour l'analyse et la visualisation de graphes.
    \item \textbf{Ollama :} pour l'utilisation locale de LLM.
\end{itemize}

\section{Matériel}

Les expériences sont menées sur un ordinateur portable équipé d'accélération GPU :
\begin{itemize}
    \item \textbf{GPU :} NVIDIA RTX 4070 Laptop (8~Go VRAM) -- pour l'entraînement et l'inférence des modèles.
    \item \textbf{CPU :} AMD Ryzen 9 8945HS (4~GHz) -- pour le prétraitement des données et l'analyse de graphes.
    \item \textbf{RAM :} 32~Go -- pour le stockage en mémoire des grandes structures de graphe.
\end{itemize}